\section{Defense Mechanisms and Execution Costs}\label{sec:defenses}
Defense placement determines both the signal changed and the party that must enforce the change.
\subsection{What a defense changes}
Table~\ref{tab:defense-space} organizes six intervention choices by enforcement point, trust, cost and residual signal. Association defenses hide a destination or its link to a user. Padding changes sizes, while scheduling changes temporal structure. Transport cover can terminate before business execution; execution cover adds real work. State isolation limits shared-resource signals, and content minimization limits what a recipient learns. These choices can be combined, but their costs and remaining views must be evaluated together.

\begin{table*}[t]\centering\small
\caption{Defense design space: enforcement authority, costs to measure and remaining exposure. Source-specific measurements appear in the text and Supplement S7.}
\label{tab:defense-space}
\setlength{\tabcolsep}{4pt}
\begin{tabularx}{\textwidth}{@{}P{.17\textwidth}P{.25\textwidth}P{.23\textwidth}Y@{}}\toprule
Layer / mechanism & What must be trusted or controlled & Costs to measure together & Residual exposure and practical decision \\\midrule
Naming / association: encrypted name resolution, relays, OHTTP & Name resolver or relay placement; the parties able to link the user to the destination & Added hops, bytes, latency and availability & Shape can survive name hiding. Choose this layer when destination association is the private target; then test remaining observers. \\\addlinespace[3pt]
Representation: message, TLS-record or packet padding & Component that controls the final visible size; padding survives later encoding & Added bytes, record/packet count, fragmentation & Timing, count and other links remain. State the padding layer and test downstream captures. \\\addlinespace[3pt]
Scheduling: batching, pacing, fixed-rate cover & Sender/relay controls output schedule and cover lifetime & First-response delay, completion delay, bandwidth during idle periods & Admission, termination and overflow can reveal activity. Budget these events as part of the visible trace. \\\addlinespace[3pt]
Transport cover: dummy exchanges discarded at a tunnel endpoint & Trusted cover terminator; no dispatch into business tools & Cover bandwidth, queueing, resource use & Other endpoints still see their inputs. Prefer this boundary when the aim is traffic shaping without extra business work. \\\addlinespace[3pt]
Execution cover: extra research or tool calls & Tool privileges, isolated state and whether decoy feedback reaches the planner & Model tokens, tool quota, task success, state changes and latency & A decoy can affect the real execution. Measure privacy and original-task utility on the same protected runs. \\\addlinespace[3pt]
Shared-state isolation / minimization & Tenant separation, cache keys and information passed between services & Lost cache reuse, compute, latency, task accuracy & Isolation removes a sharing channel, not every traffic signal. Minimize tool arguments and feedback according to each recipient's role. \\\bottomrule
\end{tabularx}
\end{table*}



\subsection{Relating communication defenses to Agent execution}
STP illustrates transport cover that can terminate at a tunnel endpoint~\cite{apthorpe2019stp}; Walkie-Talkie and FRONT provide shaping mechanisms~\cite{wang2017walkietalkie,gong2020front}. For the ticket Agent, this could mean padding the tool response at an egress proxy, pacing model tokens, or sending cover to a trusted terminator. Calling the refund tool as a decoy would instead cross the business-state boundary. The intervention must therefore specify both where cover is removed and what work it can trigger. Constant-rate shaping in the style of BuFLO and Tamaraw~\cite{dyer2012peekaboo,cai2014tamaraw} fits Agent traffic poorly. Model waits create long idle periods that the defense must fill, and streamed responses can exceed any fixed schedule. In our pilot, joint model-link shaping cost 19.5--22.5$\times$ bytes, and 1,818 of 2,736 shaped responses overran the minimum schedule, with no detectable additional reduction in recognition or localization (Section~\ref{sec:live-study}).

The measured design points occupy different parts of this space. WRA changes executed research, live OHTTP changes destination association, and the paired forwarding study changes the context sent to the model. Section~\ref{sec:evidence} gives the source-specific recovery and cost results; Section~\ref{sec:experiment} tests cover execution, forwarding and connection policies. Supplement S7 retains the complete defense map, including identity obfuscation and adversarial-Agent blocking as distinct objectives.

\subsection{Choosing a layer in practice}
Start from the private object and the observer's rights. Protect destination associations when contacted domains expose the query; protect visible model output sizes and timing when those carry the target; minimize sensitive arguments when the recipient is the observer. Then test what remains on other permitted links. A shared gateway changes which tasks contribute to a connection, so its connection policy belongs in the evaluation. For executed cover, bind privacy, token/tool use, requested-task completion and whole-state correctness to the same protected runs. Where the observer sees only padded metadata, terminate cover at a proxy and spend the budget on schedule coverage; execution depth added no measured privacy (Section~\ref{sec:cover-study}). These choices turn the risk map into a deployment-specific protection plan.
